\documentclass{article}

\usepackage[english]{babel}
\usepackage[letterpaper,top=2cm,bottom=2cm,left=3cm,right=3cm,marginparwidth=1.75cm]{geometry}
\setlength{\parindent}{0pt}

\usepackage{amsmath}
\usepackage{amsthm}
\usepackage{amssymb}
\usepackage{graphicx}
\usepackage{float}
\usepackage{listings}
\usepackage[colorlinks=true, allcolors=blue]{hyperref}
\usepackage{xcolor}

\lstset{
  basicstyle=\ttfamily\small,
  numbers=left,
  numberstyle=\tiny,
  stepnumber=1,
  numbersep=8pt,
  backgroundcolor=\color{gray!10},
  frame=single,
  rulecolor=\color{gray!60},
  tabsize=2,
  captionpos=b,
  breaklines=true,
  breakatwhitespace=true,
  showstringspaces=false,
  keywordstyle=\color{blue},
  commentstyle=\color{green!50!black},
  stringstyle=\color{orange},
}

\DeclareMathOperator{\VC}{VCdim}
\DeclareMathOperator{\KL}{KL}
\DeclareMathOperator{\supp}{supp}
\DeclareMathOperator{\sign}{sign}
\DeclareMathOperator*{\argmin}{arg\,min}

\newcommand{\EE}{\mathbb{E}}
\newcommand{\PP}{\mathbb{P}}
\newcommand{\cH}{\mathcal{H}}
\newcommand{\cD}{\mathcal{D}}
\newcommand{\cX}{\mathcal{X}}

\title{DSC 291: Learning Theory Assignment 4}
\author{Tongle Shen}

\begin{document}
\maketitle

\section{Part A: Sparse Linear Predictors and Model Selection}

\subsection{Problem 1. VC Dimension Through Supports}

For each support set $I \subseteq [d]$ with $|I|=k$, define
$$
\cH_I
=
\left\{
x \mapsto \sign(\langle w,x\rangle)
:
\supp(w)\subseteq I
\right\}.
$$
Then every $k$-sparse predictor belongs to $\cH_I$ for some $I$ of size $k$, so
$$
\cH_k = \bigcup_{|I|=k} \cH_I.
$$

Fix such an $I$. After restricting to the coordinates in $I$, the class $\cH_I$ is just the class of homogeneous halfspaces in $\mathbb{R}^k$, so
$$
\VC(\cH_I) \le k.
$$
Hence, for every $n\ge k$, Sauer--Shelah gives
$$
\Gamma_{\cH_I}(n)
\le
\sum_{j=0}^k \binom{n}{j}
\le
\left(\frac{en}{k}\right)^k.
$$

There are $\binom{d}{k}$ possible supports, so
$$
\Gamma_{\cH_k}(n)
\le
\binom{d}{k}\left(\frac{en}{k}\right)^k
\le
\left(\frac{ed}{k}\right)^k \left(\frac{en}{k}\right)^k.
$$

Now let $m=\VC(\cH_k)$. If $\cH_k$ shatters $m$ points, then
$$
2^m \le \Gamma_{\cH_k}(m).
$$
Also $\cH_k$ is a subclass of all homogeneous halfspaces in $\mathbb{R}^d$, and that larger class has VC dimension $d$, so $m\le d$. Therefore
$$
2^m
\le
\left(\frac{ed}{k}\right)^k \left(\frac{em}{k}\right)^k
\le
\left(\frac{ed}{k}\right)^{2k}.
$$
Taking logarithms,
$$
m\log 2 \le 2k \log\!\left(\frac{ed}{k}\right).
$$
So
$$
\VC(\cH_k)
\le
\frac{2}{\log 2}\,k\log\!\left(\frac{ed}{k}\right)
=
O\!\left(k\log\!\left(\frac{ed}{k}\right)\right).
$$

This is the desired upper bound.

\subsection{Problem 2. Penalty-Based SRM}

Choose, for example,
$$
p_k = \frac{6}{\pi^2 k^2},
\qquad
k=1,\dots,d.
$$
Then $\sum_{k=1}^d p_k \le 1$.

Let $C_{\mathrm{SRM}}$ be the universal constant from the class-level SRM theorem in the Week 4 lecture. That theorem says that with probability at least $1-\delta$, simultaneously for every $k$ and every $h\in\cH_k$,
$$
\left|L_{\cD}(h)-L_S(h)\right|
\le
C_{\mathrm{SRM}}
\sqrt{
\frac{
\VC(\cH_k)+\log(1/p_k)+\log(1/\delta)
}{n}
}.
$$
Using Problem 1 and enlarging the constant if necessary, there is a universal $C>0$ such that on the same event,
$$
\left|L_{\cD}(h)-L_S(h)\right|
\le
C
\sqrt{
\frac{
k\log(ed/k)+\log(1/p_k)+\log(1/\delta)
}{n}
}.
$$

Define
$$
\operatorname{pen}(k,n,\delta)
:=
C
\sqrt{
\frac{
k\log(ed/k)+\log(1/p_k)+\log(1/\delta)
}{n}
}
$$
and choose
$$
\widehat{h}
\in
\argmin_{1\le k\le d,\ h\in\cH_k}
\left\{
L_S(h)+\operatorname{pen}(k,n,\delta)
\right\}.
$$

On the high-probability event above, for every $k$ and every $h\in\cH_k$,
\begin{align*}
L_{\cD}(\widehat{h})
&\le
L_S(\widehat{h})+\operatorname{pen}(\widehat{k},n,\delta) \\
&\le
L_S(h)+\operatorname{pen}(k,n,\delta) \\
&\le
L_{\cD}(h)+2\operatorname{pen}(k,n,\delta).
\end{align*}
Taking the infimum over all $k$ and $h\in\cH_k$ gives
$$
L_{\cD}(\widehat{h})
\le
\inf_{1\le k\le d,\ h\in\cH_k}
\left\{
L_{\cD}(h)+2\operatorname{pen}(k,n,\delta)
\right\}.
$$
Absorbing the factor $2$ into the universal constant, we obtain the oracle inequality
$$
L_{\cD}(\widehat{h})
\le
\inf_{1\le k\le d,\ h\in\cH_k}
\left\{
L_{\cD}(h)
+
C
\sqrt{
\frac{
k\log(ed/k)+\log(1/p_k)+\log(1/\delta)
}{n}
}
\right\}.
$$

\subsubsection{Interpretation of the Penalty}

\textbf{Support-selection cost.}
The term $k\log(ed/k)$ is the within-class statistical cost of searching over a $k$-dimensional separator and over the $\binom{d}{k}$ possible coordinate supports. This is the price of not knowing in advance which features matter.

\textbf{Sparsity-level cost.}
The term $\log(1/p_k)$ is the across-class cost for allowing the learner to choose the complexity level $k$ after seeing the data. If we take $p_k \propto 1/k^2$, then this cost is only $O(\log k)$.

\subsection{Problem 3. Comparison with Dense Halfspaces}

Assume there exists $w^* \in \mathbb{R}^d$ with $\|w^*\|_0 \le s$ and
$$
L_{\cD}(w^*) \le \eta.
$$
Applying the oracle inequality from Problem 2 with $k=s$ and $h=w^*$ yields, with probability at least $1-\delta$,
$$
L_{\cD}(\widehat{h})
\le
\eta
+
C
\sqrt{
\frac{
s\log(ed/s)+\log(1/p_s)+\log(1/\delta)
}{n}
}.
$$

Therefore it is sufficient to take
$$
n
\ge
C'
\frac{
s\log(ed/s)+\log(1/p_s)+\log(1/\delta)
}{
\varepsilon^2
}
$$
for a universal constant $C'$, because then
$$
L_{\cD}(\widehat{h}) \le \eta+\varepsilon
$$
with probability at least $1-\delta$.

For comparison, if we learn over all homogeneous halfspaces in $\mathbb{R}^d$, then $\VC= d$, so the standard VC bound gives sample size
$$
n_{\mathrm{dense}}
=
O\!\left(
\frac{d+\log(1/\delta)}{\varepsilon^2}
\right).
$$

The sparse bound is smaller when
$$
s\log(ed/s)+\log(1/p_s) \ll d.
$$
With the concrete choice $p_k \propto 1/k^2$, this is essentially the regime
$$
s\log(ed/s) \ll d.
$$
So SRM helps exactly when the target is much sparser than the ambient dimension.

\subsection{Problem 4. Validation as Model Selection}

Write $m=n/2$. For each $k=1,\dots,d$, let
$$
w_k \in \argmin_{\|w\|_0\le k} L_{S_1}(w),
$$
and let
$$
\widehat{k}\in\argmin_{1\le k\le d} L_{S_2}(w_k).
$$
I denote the selected predictor by $w_{\widehat{k}}$.

\subsubsection{Step 1: Control the Predictors Trained on S1}

For each fixed $k$, $w_k$ is an ERM over $\cH_k$ based on the sample $S_1$ of size $m$. Since
$$
\VC(\cH_k)=O\!\left(k\log\!\left(\frac{ed}{k}\right)\right),
$$
the Week 3 ERM guarantee gives the following: for some universal constant $C_1$, with probability at least $1-\delta/(2d)$,
$$
L_{\cD}(w_k)
\le
\inf_{\|w\|_0\le k} L_{\cD}(w)
+
C_1
\sqrt{
\frac{
k\log(ed/k)+\log(2d/\delta)
}{m}
}.
$$
Applying a union bound over $k=1,\dots,d$, with probability at least $1-\delta/2$ we have simultaneously for all $k$:
$$
L_{\cD}(w_k)
\le
\inf_{\|w\|_0\le k} L_{\cD}(w)
+
\alpha_k,
$$
where
$$
\alpha_k
:=
C_1
\sqrt{
\frac{
k\log(ed/k)+\log(2d/\delta)
}{m}
}.
$$

\subsubsection{Step 2: Use S2 to Compare the d Candidates}

Now condition on $S_1$. Then the candidates $w_1,\dots,w_d$ are fixed predictors, and $S_2$ is an independent sample of size $m$.

By Hoeffding's inequality plus a union bound over the $d$ fixed candidates, with probability at least $1-\delta/2$,
$$
\left|L_{\cD}(w_k)-L_{S_2}(w_k)\right|
\le
\beta
\qquad
\text{for all } k=1,\dots,d,
$$
where
$$
\beta
:=
\sqrt{
\frac{
\log(4d/\delta)
}{
2m
}
}.
$$

\subsubsection{Step 3: Combine the Two Events}

Assume both high-probability events from Steps 1 and 2 hold. Then for every $k$,
\begin{align*}
L_{\cD}(w_{\widehat{k}})
&\le
L_{S_2}(w_{\widehat{k}})+\beta \\
&\le
L_{S_2}(w_k)+\beta \\
&\le
L_{\cD}(w_k)+2\beta \\
&\le
\inf_{\|w\|_0\le k} L_{\cD}(w)+\alpha_k+2\beta.
\end{align*}
Taking the infimum over $k$ gives
$$
L_{\cD}(w_{\widehat{k}})
\le
\inf_{1\le k\le d,\ \|w\|_0\le k}
\left\{
L_{\cD}(w)+\alpha_k+2\beta
\right\}.
$$

Since $m=n/2$, both $\alpha_k$ and $\beta$ are of order
$$
\sqrt{
\frac{
k\log(ed/k)+\log(d/\delta)
}{n}
},
$$
after adjusting constants. Therefore, for some universal constant $C$,
$$
L_{\cD}(w_{\widehat{k}})
\le
\inf_{1\le k\le d,\ \|w\|_0\le k}
\left\{
L_{\cD}(w)
+
C
\sqrt{
\frac{
k\log(ed/k)+\log(d/\delta)
}{n}
}
\right\}.
$$

\subsubsection{Comparison with Penalty-Based SRM}

\textbf{Penalty-based SRM.}
All $n$ samples are used both to fit predictors and to choose the sparsity level through the explicit penalty. The price for adapting to $k$ is $\log(1/p_k)$.

\textbf{Validation.}
The first half $S_1$ is used only for fitting the predictors $w_k$, while the second half $S_2$ is used only for choosing among those $d$ fitted candidates. The price for this separation is twofold:
\begin{itemize}
\item each predictor is trained on only $n/2$ samples rather than all $n$ samples;
\item model selection over the $d$ candidates contributes an extra $\log d$ term.
\end{itemize}
So validation is conceptually simple, but it is statistically less efficient than using one SRM objective on the full sample.

\subsection{Problem 5. Computation}

Now assume the realizable case: the sample is consistent with some $k$-sparse homogeneous halfspace.

\subsubsection{Brute-Force Algorithm}

Enumerate all support sets $I\subseteq [d]$ with $|I|=k$. For each such $I$, solve the feasibility problem
$$
y_i \langle w_I, x_{i,I}\rangle \ge 1
\qquad
\text{for } i=1,\dots,n,
$$
where $x_{i,I}$ is the restriction of $x_i$ to the coordinates in $I$.

This feasibility problem is valid because if the sample is realizable on support $I$, then there exists a separator with strictly positive margin on the finite sample, and we can rescale it so that all margins are at least $1$.

If the feasibility problem is solvable for some $I$, output the corresponding classifier. If no support works, declare failure.

\subsubsection{Runtime}

There are
$$
\binom{d}{k}
\le
\left(\frac{ed}{k}\right)^k
$$
possible supports. For each support we solve a linear feasibility problem with $k$ variables and $n$ constraints, which takes $\operatorname{poly}(n,k)$ time.

So the total runtime is
$$
T(d,k,n)
=
\binom{d}{k}\operatorname{poly}(n,k)
\le
\left(\frac{ed}{k}\right)^k \operatorname{poly}(n,k).
$$

\subsubsection{When Is This Polynomial in d?}

This runtime is polynomial in $d$ only when $k$ is a fixed constant. If $k$ grows with $d$, then the combinatorial factor $\binom{d}{k}$ gives a superpolynomial search cost. For example:
\begin{itemize}
\item if $k=O(1)$, the runtime is polynomial in $d$;
\item if $k=\Theta(\log d)$, the runtime is already quasipolynomial or worse;
\item if $k=\Theta(d^\alpha)$ for any fixed $\alpha>0$, the runtime is exponential in a power of $d$.
\end{itemize}

\subsubsection{Lesson}

This is exactly the Week 4 distinction between sample complexity and computational complexity. Statistically, SRM shows that sparse predictors can be learned from about
$$
O\!\left(\frac{k\log(ed/k)}{\varepsilon^2}\right)
$$
samples. Computationally, however, finding the best sparse support may require brute-force search over $\binom{d}{k}$ possibilities. So favorable sample complexity does not by itself imply an efficient learning algorithm.

\newpage
\section{Part B: PAC-Bayes for Thresholds}

Throughout this part,
$$
\cX_N=\{1,2,\dots,N\},
\qquad
\cH_N=\{h_t:t\in\{1,\dots,N+1\}\},
\qquad
h_t(x)=\mathbf{1}[x\ge t].
$$
I write the lecture PAC-Bayes complexity term as
$$
\mathrm{PB}(Q,P,n,\delta)
:=
\sqrt{
\frac{
\KL(Q\Vert P)+\log(2n/\delta)
}{
2(n-1)
}
}.
$$
Then the lecture theorem says that with probability at least $1-\delta$, simultaneously for all posteriors $Q$,
$$
L_{\cD}(Q)\le L_S(Q)+\mathrm{PB}(Q,P,n,\delta).
$$

\subsection{Problem 1. VC Dimension and Point-Posterior PAC-Bayes}

\subsubsection{VC Dimension}

First, $\VC(\cH_N)\ge 1$ because any single point $x\in\cX_N$ can be labeled both ways:
\begin{itemize}
\item $h_1(x)=1$,
\item $h_{N+1}(x)=0$.
\end{itemize}

Now take two points $x<y$. Every threshold labeling is monotone: once the label becomes $1$, it stays $1$ for all larger points. Therefore the pattern
$$
(h(x),h(y))=(1,0)
$$
is impossible. So no two-point set is shattered, and hence
$$
\VC(\cH_N)=1.
$$

\subsubsection{Point Posterior}

Let $P$ be the uniform prior on $\cH_N$, so
$$
P(h_t)=\frac{1}{N+1}
\qquad
\text{for all } t.
$$
Let $\widehat{h}_t=A(S)$ be any deterministic consistent ERM, and define
$$
Q_S=\delta_{\widehat{h}_t}.
$$
Then
$$
\KL(Q_S\Vert P)
=
\log\!\left(\frac{1}{P(\widehat{h}_t)}\right)
=
\log(N+1).
$$

Because we are in the realizable setting and $\widehat{h}_t$ is consistent,
$$
L_S(Q_S)=L_S(\widehat{h}_t)=0.
$$
Therefore PAC-Bayes gives
$$
L_{\cD}(Q_S)
=
L_{\cD}(\widehat{h}_t)
\le
\sqrt{
\frac{
\log(N+1)+\log(2n/\delta)
}{
2(n-1)
}
}
$$
with probability at least $1-\delta$.

\subsubsection{Why This Depends on log N}

This certificate pays $\log(N+1)$ because a point posterior must identify one specific threshold out of $N+1$ possibilities under the uniform prior. So PAC-Bayes is acting like a finite-class description-length bound.

By contrast, the VC guarantee for thresholds uses the order structure of the whole class and sees only $\VC(\cH_N)=1$. That structural argument is independent of $N$, so its realizable sample complexity does not grow with $\log N$.

\subsection{Problem 2. Version-Space Posterior}

Define
$$
a(S)=\max\{x_i:y_i=0\},
\qquad
b(S)=\min\{x_i:y_i=1\},
$$
with the conventions $a(S)=0$ if there are no negatives and $b(S)=N+1$ if there are no positives.

\subsubsection{Description of the Version Space}

A threshold $h_t$ is consistent with $S$ iff:
\begin{itemize}
\item every negative example satisfies $x_i<t$, so $t>a(S)$;
\item every positive example satisfies $x_i\ge t$, so $t\le b(S)$.
\end{itemize}
Hence
$$
V(S)=\{t:a(S)<t\le b(S)\}.
$$
Since $t$ is integer-valued, this set is
$$
\{a(S)+1,\dots,b(S)\},
$$
so
$$
|V(S)|=b(S)-a(S).
$$

\subsubsection{KL of the Uniform Version-Space Posterior}

Let $Q_V$ be the uniform distribution on $\{h_t:t\in V(S)\}$. Then
$$
Q_V(h_t)=\frac{1}{|V(S)|}
\qquad
\text{for } t\in V(S).
$$
Since $P$ is uniform on all $N+1$ thresholds,
\begin{align*}
\KL(Q_V\Vert P)
&=
\sum_{t\in V(S)}
\frac{1}{|V(S)|}
\log\!\left(
\frac{1/|V(S)|}{1/(N+1)}
\right) \\
&=
\log\!\left(\frac{N+1}{|V(S)|}\right).
\end{align*}

\subsubsection{Empirical Risk}

Every threshold in $V(S)$ is consistent with the sample, so each one has empirical risk $0$. Therefore
$$
L_S(Q_V)=\EE_{h\sim Q_V}[L_S(h)]=0.
$$
The PAC-Bayes bound becomes
$$
L_{\cD}(Q_V)
\le
\sqrt{
\frac{
\log((N+1)/|V(S)|)+\log(2n/\delta)
}{
2(n-1)
}
}.
$$

\subsubsection{Why It Can Be Better Than a Point Posterior}

The point posterior pays $\log(N+1)$, while the version-space posterior pays only
$$
\log\!\left(\frac{N+1}{|V(S)|}\right).
$$
So if many thresholds are consistent, spreading the posterior uniformly across them lowers the KL term by $\log|V(S)|$ without increasing empirical loss. That can make the PAC-Bayes certificate strictly tighter.

\subsection{Problem 3. Spreading Helps KL, but Can Hurt True Risk}

Let $\cD$ be the realizable distribution with uniform marginal on $\cX_N$ and labels generated by the true threshold $h_\tau$.

\subsubsection{Risk Formula}

For any threshold $h_t$, the disagreement set between $h_t$ and $h_\tau$ is exactly the integer interval between $t$ and $\tau$:
\begin{itemize}
\item if $t>\tau$, disagreement occurs on $\tau,\tau+1,\dots,t-1$;
\item if $t<\tau$, disagreement occurs on $t,t+1,\dots,\tau-1$.
\end{itemize}
In both cases the number of disagreements is $|t-\tau|$. Since the marginal on $\cX_N$ is uniform,
$$
L_{\cD}(h_t)=\frac{|t-\tau|}{N}.
$$
Therefore for any nonempty $W\subseteq\{1,\dots,N+1\}$ and uniform posterior $Q_W$,
\begin{align*}
L_{\cD}(Q_W)
&=
\frac{1}{|W|}
\sum_{t\in W} L_{\cD}(h_t) \\
&=
\frac{1}{N|W|}
\sum_{t\in W}|t-\tau|.
\end{align*}

\subsubsection{Concrete Example}

Take
$$
N=20,
\qquad
\tau=11,
\qquad
S=\{(5,0),(17,1)\}.
$$
This sample is realizable by $h_{11}$, since $5<11$ and $17\ge 11$.

Here
$$
a(S)=5,
\qquad
b(S)=17,
$$
so
$$
V(S)=\{6,7,\dots,17\}
$$
and
$$
|V(S)|=12.
$$

For the uniform prior $P$,
$$
\KL(Q_{V(S)}\Vert P)
=
\log\!\left(\frac{21}{12}\right)
<
\log(21)
=
\KL(\delta_{h_{11}}\Vert P).
$$

On the other hand,
$$
L_{\cD}(\delta_{h_{11}})=0
$$
because $h_{11}$ is the true threshold, while
\begin{align*}
L_{\cD}(Q_{V(S)})
&=
\frac{1}{20\cdot 12}
\sum_{t=6}^{17}|t-11| \\
&=
\frac{36}{240}
=
\frac{3}{20}
>
0.
\end{align*}
So this example satisfies
$$
\KL(Q_{V(S)}\Vert P) < \KL(\delta_{h_{11}}\Vert P)
$$
but
$$
L_{\cD}(Q_{V(S)}) > L_{\cD}(\delta_{h_{11}}).
$$

\subsubsection{Why This Does Not Contradict PAC-Bayes}

PAC-Bayes gives an upper bound on true risk; it does not say that a posterior with smaller KL must have smaller true risk. Lowering KL by spreading mass can also move probability onto hypotheses that make more mistakes under the data distribution.

So there is no contradiction: $Q_{V(S)}$ has a better complexity term, but it averages over many wrong thresholds, and that can increase the actual risk.

\subsection{Problem 4. Every Fixed Prior Can Be Attacked}

Let $P$ be any prior on $\cH_N$, fixed before the sample, and assume $N\ge 3$.

\subsubsection{Choose a Low-Mass Internal Threshold}

There are $N-1$ internal thresholds $\{h_2,\dots,h_N\}$. Their total prior mass is at most $1$, so by averaging there exists some
$$
\tau\in\{2,\dots,N\}
$$
such that
$$
P(h_\tau)\le \frac{1}{N-1}.
$$

\subsubsection{A Realizable Adversarial Distribution}

Define $\cD_\tau$ by
$$
\PP_{(x,y)\sim\cD_\tau}[x=\tau-1,y=0]=\frac12,
\qquad
\PP_{(x,y)\sim\cD_\tau}[x=\tau,y=1]=\frac12.
$$
This distribution is realizable by $h_\tau$.

The sample fails to contain both support points only if all $n$ observations are $(\tau-1,0)$ or all $n$ observations are $(\tau,1)$. Therefore
$$
\PP_{S\sim \cD_\tau^n}[\text{$S$ contains both support points}]
=
1-2\left(\frac12\right)^n
=
1-2^{1-n}.
$$

\subsubsection{The Version Space Collapses}

Assume the event above holds. Then $S$ contains at least one negative example at $\tau-1$ and at least one positive example at $\tau$. So any consistent threshold must satisfy
$$
t>\tau-1
\qquad
\text{and}
\qquad
t\le \tau.
$$
Since $t$ is an integer, the only possibility is
$$
t=\tau.
$$
Hence
$$
V(S)=\{\tau\}.
$$

\subsubsection{Any Zero-Empirical-Error Posterior Must Be a Point Mass}

Let $Q$ be any posterior with $L_S(Q)=0$. Every threshold $h_t$ with $t\ne\tau$ misclassifies at least one of the two support points $(\tau-1,0)$ and $(\tau,1)$, and both points appear in $S$. Therefore every such $h_t$ has strictly positive empirical risk on $S$.

Since
$$
L_S(Q)=\EE_{h\sim Q}[L_S(h)]
$$
is an average of nonnegative numbers and equals $0$, $Q$ can place mass only on zero-empirical-error hypotheses. But the only such hypothesis is $h_\tau$. Therefore
$$
Q=\delta_{h_\tau}.
$$

Consequently,
$$
\KL(Q\Vert P)
=
\log\!\left(\frac{1}{P(h_\tau)}\right)
\ge
\log(N-1).
$$

\subsubsection{What This Does and Does Not Show}

This lower bound shows a limitation of zero-empirical-error PAC-Bayes certificates with a fixed prior: for some realizable threshold problems, every such certificate must pay at least $\log(N-1)$ in complexity.

It does \emph{not} show that thresholds are statistically hard to learn. Thresholds still have VC dimension $1$, so their PAC sample complexity is independent of $N$. The lower bound is about one specific style of PAC-Bayes certificate, not about the intrinsic learnability of the threshold class.

\newpage
\section{Part C: AI Usage Report}

To be honest, I'm so busy this week so here we keep the original skill and workflow as previous assignments.
But when I'm writing research papers, I have tried to build a whole workflow to build the theorms and proof chain.
I found it not as easy as this homework solver, because the corpus is noisy and actually the codex is too conservative to the new
or frontier concepts. Also, the requirements for a research paper is much more strict, and you need to be very careful with the whole logic chain.

Anyway, I still think that we can build a workflow to write the research paper, but we need much more time for verification and modifying the skills.

\end{document}
